Requirements of Approach Section (2000 words maximum)

Explain how you have designed your work so that it:
is effective and appropriate to achieve your objectives
is feasible, and comprehensively identifies any risks to delivery and how you will manage them
uses a clearly written and transparent methodology (if applicable)
summarises the previous work and describes how you will build on and progress this work (if applicable)
will maximise translation of outputs into outcomes and impacts
Within the Approach section we also expect you to:
demonstrate access to the appropriate services, facilities, infrastructure, or equipment to deliver the proposed work
provide a detailed and comprehensive project plan, including milestones and timelines in the form of a Gantt chart or diagram
explain who you intend to collaborate with at the host organisation and your plans for wider research collaborations (project partner details should be provided only in that section)
include details of work that will take place as part of the proposed fellowship at a second UK or overseas organisation (if applicable)
explain and justify how you will approach diversity and inclusion in the study population and follow the MRC embedding diversity in research design policy (if applicable)
show how you will use male and female animals or tissues and cells from female and male donors (if applicable) in your research. If you are not proposing to do this, justify why
explain and justify the inclusion of public partnerships (if applicable) and the added value these offence

Including ‘References’ and ‘Images’
You are permitted to include ‘References’ and ‘Images’ within this section. If you choose to include these, you must refer to the guidance within the ‘Vision’ section. Failure to follow the guidance may result in your application being rejected.
You should detail any information regarding your statistical design in the 'Reproducibility and statistical design' section.


_______________________________

\section*{3. Approach}
\subsection*{Aims}


\subsection*{Background}


\subsection*{Research Strategy}


\subsection*{WP1. Horizontal Gene Transfer of Klebsiella species, a Global AMR Pathogen}
The Floto Lab is conducting an in depth study of Klebsiella. Metadata I have curated including clinical outcomes of focus as well as leading the work on study of mobile genetic elements. Pictures of the dataset (global distribution and phylogeny). 

a)  Mobile Genetic Elements in Klebsiella 

b)  Role and Diversity of Plasmids in Klebsiella

This is clearly multifaceted and involves evolution of both the core genome, vertical transmission, phylogeny, pathoadaption, dating, variant calling and hotspot analysis as well as the horizontal gene transmission that I will focus on.  It requires a team and I will be working alongside Weimann, Dinan, Figuorea, Ruiz, and Bruchmann. The vast experience using the bioinformatic tools involved and collaboration of this team of bioinformaticians derisks the output of my work while I focus on the particular aspects of mobile elements and plasmids in the overall pathoadaption. 

\subsection*{WP2. Predict Behaviour of Klebsiella and Pseudomonas Aeruginosa from Genomes}
Mini-background - My plots for Bacformer for prediction and gene association using explainability and attribution with Captum via Input Gradient and Gene Ablation and UMAP of  the protein embedding hotspots. 
My plots for XGBoost predictions of Klebseilla and PsA across antibiotics

a)  AMR - Predictions, generalisation to larger set 
a)  Predicting AMR in PsA from WGS in Klebsiella and PsA
Dataset - EBI AST results summarise
i)   Test accuracy of Bacformer 
ii)  Accuracy of a one hot encoded labels for AMR associated genes (gene presence / absence from AMR Finder Plus through NCBI) -->  put through an XGBoost Model
iii)  If Can we improve on Bacformer by incorporating both into embeddings - label each gene by AMR Finder plus result for antibiotic of relevance and read in whole sequence
Given that AMR is well predicted in Klebsiella by screening on smaller NCBI dataset with XGBoost model where we already achieved AUROC of 0.8-0.99 across antibiotics - Figure to go in Figures - for Klebsiella and Pseudomonas, 
Assess ability to remove AST testing from a bioinformatic workflow where WGS is already being done, informing our testing in part 3. 

b)  Invasive Disease - 
Dataset 
Other invasive phenotypes
Predictions of Bacformer vs pyseer and correlations with plasmid analysis / Kleborate output / hotspot analysis. 
Uncovering the basis of resistance with Pyseer from predictions on anything that is initially predicted poorly


\subsection*{WP3.  Clinical antimicrobial resistance and resilience}
Background: Fully collected and funded clinical study for which I am PI and have been working on for a number of years with CCLI. More than 600 courses of intravenous antibiotics including over 1200 antibiotics, 1200 sputum samples and 1200 endpoints. 

Background of the Kishony paper demonstrating strain (and species) switching underlying gain in resistance after antibiotic treatment and treatment failure using SNP distances. 

The main risk here is that in respiratory disease the consistency of the culture and AST results within individuals over time in sputum samples (reference) despite variable response to treatment over time for the same individual (reference, or quote CCLI / accepted experience - quote from microbial guidelines needed). But this belies the polyclonality and diverse microbial species that make up the respiratory sputum, often signature to patients over time without change through antibiotic courses (reference (Figure)

a)  Treatment efficacy per culture Antibiotic Sensitivity Test (AST)
Using historic culture results available at the time of consultation (when starting or changing antibiotics) and focussed on Pseudomonas Aeruginosa with secondary results for all other non-mycobacterial cultures, calculate efficacy of treatment by:
i)  Global change scale
ii)  Changes in symptom questionnaire
iii)  Changes in Chronic Airway Assessment Test (CAAT)
iv)  Time till next treatment (stratified into < 3 months, < 12 months or > 12 months

Correlate these results to AST results with subgroup analysis by underlying disease and species on culture (for non-pseudomonas cultures).  

b)  Strain Identification and Switching to Predict Outcomes
With microbiologists in the Floto lab and at the Bryant lab (Sanger) we will culture sputum and conduct AST testing. 
We will do WGS of strains to identify strain switching, accessory genome and plasmid in the genomes. 
The polymicrobial and polystrain/polyclonal structure of bacterial communities can be lost in cultures and thus in WGS.  

This dense dataset will allow iterative testing of predictions but needs careful development of how genomics can be used to predict outcomes.
our first focus will be on streain switching as a cause of treatment failure and resistance emergence, given the Kishony paper results. 
In future we may look at bacteria killing dynamics (ref. Boeck and TB paper on antimicrobial resilience - Liu and Fortune)

We want an iterative approach to affordably test hypotheses and approaches. 
Our initial plan is to focus on infection with pseudomonas aeruginosa as a model system for respiratory disease and infection. 
i/  We will screen patients based on previous clinical cultures of pseudomonas aeruginosa
ii/  We will conduct 16S RNA guides of whole sputum samples to characterise the diversity and abundance at the species level and confirm PsA presence
iii/  We will then culture of PsA specific plates
-  We will separate colonies on the plate
-  Different colonies will have PCR guide testing for strain lineage switching developed by the Bryant lab
-  Distinct lineages will be selected for WGS
-  A separate sweep metagenome WGS will also be conducted
-  AST results will, in first instance, be predicted from genomic features, leveraging our work in section 2a with the knowledge that our current XGBoost model already gives excellent AUROC and AUPRC as baseline

We will then study strain switching by SNP distance and check if AST of the different strains predicts treatment failure and resistance directly.  
We will also consider using Bacformer to predict clinical outcome directly from the genome, rather than with AST results as a proxy. 


i)  We will use RNA guides (developed by the Bryant lab) to track specific accessories genes and label specific sublineages within the plate
ii)  AST and WGS of multiple colonies on the plate to test for polyclonality and variable-AST response within the culture plate.  We will do this selectively using RNA guides to look at core genes in the cultures species to look for different phylogenies which we will select for WGS
iii)  16S RNA guides of whole sputum samples to characterise the diversity and abundance at the species level

This data is the minimum the might predict strain switching in respiratory samples.

ML model of genomic data as well as clinical data (culture history and previous WGS data + BSI score, medical history) to predict efficiency and resistance to antibiotics (c.f Stracey). 






\subsection*{Outcomes for the Field and Career}
i)  Generate models of Klebsiella pathogen evasion and evolution at global population level and in clinical settings using cutting edge bioinformatics and machine learning. 
ii)  Position me as a leader in the field of Klebsiella pathogen evasion and evolution using cutting edge bioinformatics and machine learning able to transfer and design studies across machine learning and into clcinical microbiocial practice
iii)  Add to the position the Floto and University of Cambridge as a leader in these fields and with a pipeline to continue this into intellectual property and commercialisation of important discoveries


\subsection*{Timeline:  See text and GANTT Chart}

\subsection*{Floto Lab Environment}
Weimann (publications)
Wendy Figueroa (publications)
Dynan (publications)
Chris Ruiz
Sebastian Bruchmann (Klebsiella publications)

\subsection*{Collaborators (Partners)}
Josie Bryant
John Lees
Adrian Cazalles - doesnt have own lab - but I have met and we will iterate with him
Mentoring by Ewan Harrison

\subsection*{Data Dependencies}
Data sources from EBI / NCBI and Bakrep
Clinical data from SputaGen study

\subsection*{Services and Equipment Required}
CPU time
GPU time
Cultures
Assembly
PhD fees for UK student

\subsection*{Risks:  Addressed above}

\subsection*{References:}
To be completed

\subsection*{Figures:} 
i)  Global distribution of the Klebsiella set
ii)  Clonal and sublineage variance of Klebsiella genomes
iii)  Curated isolation source and host species
iv)  Pathoadaption of PsA (from Weimann et al)
vii)  Dating phylogenies in PsA 
v)  Global spread of M. abscessus (from Bryant et al)
vi)  Identifying genes with GWAS from Pyseer and Hotspot analysis
vii)  Bacformer architecture and predictions - AUROC and AUPRC, plus identifying genes using explainability techniques from Captum
viii)  Exploring the embedding space of Bacformer
ix)  Clustering of the pangenome with Panaroo
x)  Trajectories of pangenome in S. Aureus
xi)  Identifying operons with Panstripe and GGCaler
xii) Sputum microbial communities in bronchiectasis - Variability of species composition of sputum between individuals, but consistent signature over time within individuals

GANTT Chart


-----------------------
